\documentclass[11pt]{article}

\usepackage[margin=1in]{geometry}
\usepackage{amsmath,amssymb}
\usepackage{booktabs}
\usepackage{longtable}
\usepackage{array}
\usepackage{enumitem}
\usepackage{xcolor}
\usepackage{hyperref}
\usepackage{microtype}
\usepackage[T1]{fontenc}
\usepackage{lmodern}

\hypersetup{
  colorlinks=true,
  linkcolor=blue!50!black,
  urlcolor=blue!50!black,
  citecolor=blue!50!black
}

\newcommand{\code}[1]{\texttt{#1}}
\newcommand{\Rpath}[1]{\path{#1}}
\newcommand{\dd}{\,\mathrm{d}}
\newcommand{\diag}{\operatorname{diag}}
\newcommand{\tr}{\operatorname{tr}}
\newcommand{\argmin}{\operatorname*{arg\,min}}
\newcommand{\reportbuilddatetime}{2026-05-15 01:56:59 EDT}

\title{gflow Graph Low-Pass GCV Operator Used by \code{fit.rdgraph.regression()} and \code{refit.rdgraph.regression()}}
\author{SIMODS manuscript provenance note}
\date{Built \reportbuilddatetime}

\begin{document}
\maketitle

\section{Purpose and Manuscript Role}

This note describes the graph low-pass smoothing operator implemented in
\code{gflow} by \code{fit.rdgraph.regression()} and reused by
\code{refit.rdgraph.regression()}. In the SIMODS non-oracle graph-selection
benchmark, this operator is intended to serve as one candidate selector: for
each candidate graph, fit a graph-spectral smoother, record its generalized
cross-validation (GCV) score, and compare candidate graphs by aggregating those
scores over selected feature signals.

Operationally, the benchmark should construct the graph/operator with
\code{fit.rdgraph.regression()}. When the same fitted graph geometry and
spectral filter are to be applied to many response or feature columns, the
benchmark should call \code{refit.rdgraph.regression()} on the fitted object.
The refit path reuses the cached graph geometry, eigenvectors, eigenvalues, and
filter family. Depending on \code{per.column.gcv}, it either reuses the
original smoothing parameter and filtered eigenvalues stored in the fitted model
or selects a feature-specific smoothing parameter in the cached eigenbasis.
Either way, it avoids rebuilding the Riemannian graph geometry for every
feature.

\paragraph{Current implementation convention confirmed from source.}
For the current \code{fit.rdgraph.regression()} low-pass operator path with a
precomputed graph supplied through \code{adj.list} and \code{weight.list},
\code{weight.list} is interpreted as a positive edge-length list used for local
neighborhood ordering and truncation, not as an affinity or conductance list.
The implementation does \emph{not} directly use
\[
  a_{ij} = \frac{1}{\ell_{ij}+\epsilon}
\]
as the smoothing edge weight. Instead, the supplied edge lengths
\(\ell_{ij}\) are used to order local graph neighborhoods. The smoothing
operator is then built from neighborhood-intersection densities computed by the
Riemannian-complex initialization; the conductance used in the
mass-symmetrized spectral Laplacian is
\[
  c_e = \frac{1}{\max(\rho_1(e),10^{-10})},
\]
where \(\rho_1(e)\) is the edge mass/density computed from neighborhood
intersections. This is the key convention for the manuscript and benchmark.

\section{Graph Input and Weight Semantics}

Let \(G=(V,E)\) be an undirected connected graph on \(n=|V|\) vertices. In the
R API, a precomputed graph is supplied through:
\begin{itemize}[leftmargin=1.5em]
  \item \code{adj.list}: a length-\(n\) list whose \(i\)th entry gives the
  1-based neighbors of vertex \(i\);
  \item \code{weight.list}: a length-\(n\) list parallel to \code{adj.list}
  whose entries are positive edge lengths \(\ell_{ij}\).
\end{itemize}

The R wrapper validates that \code{adj.list} and \code{weight.list} are supplied
together, have matching lengths, contain no self-loops or duplicate neighbors,
are undirected, and have reciprocal edge weights agreeing within tolerance. It
then converts the adjacency list to 0-based indices for C++ while passing the
numeric weights through as doubles.

The supported R-level precomputed-graph API currently requires
\code{use.counting.measure = TRUE}. The wrapper explicitly stops if a
precomputed \code{adj.list}/\code{weight.list} is supplied with
\code{use.counting.measure = FALSE}. Therefore, in the benchmark usage
envisioned here, \(\ell_{ij}\) does not define a distance-based vertex reference
measure; it defines edge lengths used to sort local neighborhoods.

\subsection{From Edge Lengths to Local Neighborhoods}

In the precomputed-graph C++ path, each row \(i\) is converted into a sorted
list of pairs \((\ell_{ij},j)\), with a self-entry \((0,i)\) inserted. The
implementation keeps the first \(\min(k_{\mathrm{cpp}}, \deg(i)+1)\) entries
as the local neighborhood \(\widehat N_k(i)\). Because the R wrapper passes
\code{k + 1L} into C++ to account for the self neighbor convention, a user-level
\code{k=K} corresponds to self plus up to \(K\) shortest graph-neighbors when
the supplied degree is at least \(K\).

Thus edge lengths enter the operator through neighborhood ordering and
truncation:
\[
  \widehat N_K(i)
  =
  \{i\}\cup
  \text{the up to \(K\) adjacent vertices with smallest } \ell_{ij}.
\]
The graph's full edge registry still contains all supplied undirected edges;
the local neighborhoods are used to compute density and intersection masses.

\subsection{Recent GCV Graph-Selection Regime}

In the current supported precomputed-graph path, and in the recommended SIMODS
graph-selection benchmark regime, calls to \code{fit.rdgraph.regression()} use
the default counting-measure regime:
\begin{itemize}[leftmargin=1.5em]
  \item \code{use.counting.measure = TRUE};
  \item \code{response.penalty.exp = 0};
  \item \code{density.alpha = 1.5};
  \item \code{density.epsilon = 1e-10}.
\end{itemize}
The first and fourth values are especially easy to misread. In this precomputed
graph path, \code{density.epsilon} is not an \(\ell_{ij}\)-to-conductance
regularizer. It belongs to the optional distance-reference branch
\((\epsilon+d_k)^{-\alpha}\), which is not used when
\code{use.counting.measure = TRUE}. The setting
\code{response.penalty.exp = 0} disables response-coherence modulation, so the
operator is the initialization-derived graph smoother rather than an
iteratively response-adapted smoother.

Some historical or pre-run scripts experimented with
\code{use.counting.measure = FALSE}. Current \code{gflow} rejects that setting
for precomputed \code{adj.list}/\code{weight.list} graphs, so those historical
calls should not be used as canonical method descriptions for the current
precomputed-graph operator.

In this default counting-measure regime, for a graph edge \(e=(i,j)\):
\begin{enumerate}[leftmargin=1.5em]
  \item The local neighborhoods \(\widehat N_k(i)\) and \(\widehat N_k(j)\)
  are built from graph neighbors sorted by the supplied positive edge lengths,
  with the vertex itself included as a zero-distance self entry.
  \item Each vertex has reference mass one, up to the global normalization
  convention for the reference measure.
  \item The raw edge mass is essentially the local-neighborhood overlap,
  \[
    \rho_{1,\mathrm{raw}}(i,j)
      =
      \left|\widehat N_k(i)\cap \widehat N_k(j)\right|.
  \]
  More generally, the code computes
  \[
    \rho_{1,\mathrm{raw}}(i,j)
      =
      \sum_{v\in \widehat N_k(i)\cap \widehat N_k(j)} \mu(v),
  \]
  and \(\mu(v)=1\) in the counting-measure case.
  \item The implementation then globally normalizes edge masses to the number
  of undirected edges, applies density floors/clamps for numerical stability,
  and stores the resulting symmetric edge density as \(\rho_1(e)\).
  \item The conductance used by the mass-symmetrized spectral operator is the
  inverse of this normalized/floored edge mass:
  \[
    c_e^\rho
      =
      \frac{1}{\max(\rho_1(e),10^{-10})}.
  \]
\end{enumerate}
This description should be read as specific to the default
counting-measure/precomputed-graph regime. Other parameter settings can change
the reference measure or density construction, and response modulation can
change edge masses during fitting.

\section{Operator Construction}

\subsection{Reference Measure and Densities}

For the supported precomputed-graph path, the reference measure is the counting
measure:
\[
  \mu(i) = 1
\]
up to the implementation's normalization convention. The C++ implementation
also contains a distance-based reference-measure branch,
\[
  \mu(i) \propto (\epsilon + d_k(i))^{-\alpha},
\]
with optional clamping and dynamic-range control, but this branch is not
available through the R precomputed-graph API because the wrapper requires
\code{use.counting.measure = TRUE} when \code{adj.list}/\code{weight.list} are
provided.

The initial vertex and edge densities are computed from local neighborhoods
and local-neighborhood intersections:
\[
  \rho_0(i)
    =
    \sum_{v\in \widehat N_K(i)} \mu(v),
  \qquad
  \rho_{1,\mathrm{raw}}([i,j])
    =
    \sum_{v\in \widehat N_K(i)\cap \widehat N_K(j)} \mu(v).
\]
The code normalizes vertex densities to sum to the number of vertices and edge
densities to sum to the number of undirected edges, applies floors/clamps for
numerical stability, and stores the resulting \(\rho_1\) edge densities
symmetrically for both endpoints.

\subsection{Mass Matrices}

The vertex mass matrix is diagonal:
\[
  M_0 = \diag(m_1,\ldots,m_n),
  \qquad
  m_i = \max(\rho_0(i),10^{-8}).
\]

The edge mass matrix \(M_1\) is populated from edge densities. In the
diagonal-only case, which is the default whenever triangle construction is not
active, the diagonal entry for edge \(e=[i,j]\) is
\[
  M_1(e,e) = \max(\rho_1(e),10^{-15}).
\]
When triangle cofaces are built, the implementation may also add off-diagonal
entries to \(M_1\) using triple neighborhood intersections:
\[
  M_1(e_{ij},e_{is})
    =
    \sum_{v\in
      \widehat N_K(i)\cap \widehat N_K(j)\cap \widehat N_K(s)}
      \rho_0(v),
\]
for edge pairs incident in a triangle. If triangle cofaces are absent, too
large, or inconsistent, the code falls back to a diagonal edge mass
approximation.

\subsection{Boundary Operator and Hodge Laplacian}

Let \(B_1\in\mathbb{R}^{n\times |E|}\) be the oriented vertex-edge incidence
matrix. For an edge \(e=(i,j)\) with \(i<j\),
\[
  B_1(i,e)=-1,\qquad B_1(j,e)=+1.
\]
The general Hodge 0-Laplacian assembled by \code{laplacian\_bundle\_t} has the
up-term
\[
  L_0 = B_1 M_1^{-1} B_1^\top M_0,
\]
with the inverse metric applied by helper routines. This object is retained in
the implementation as \code{L[0]}.

\subsection{Actual Spectral Operator Used for Smoothing}

The spectral decomposition used for smoothing preferentially uses a
mass-symmetrized vertex Laplacian \(\widetilde L_0\), named
\code{L0\_mass\_sym} in the implementation:
\[
  \widetilde L_0
    =
    M_0^{-1/2}\, L_{\mathrm{div}}\, M_0^{-1/2},
  \qquad
  L_{\mathrm{div}} = B_1 C B_1^\top.
\]
Here \(C=\diag(c_e)\) is the diagonal conductance matrix and
\[
  c_e = \frac{1}{\max(\rho_1(e),10^{-10})}.
\]
The code comments describe this as the DEC convention \(c_e=1/m_1(e)\), so a
larger edge mass implies weaker penalization of variation across that edge. If
\code{L0\_mass\_sym} is unavailable, the eigensolver falls back to \code{L[0]},
but the normal path is the mass-symmetrized operator.

This distinction matters for manuscript wording. A simple graph-Laplacian
intuition is acceptable if stated as an intuition:
shorter supplied edge lengths affect which vertices enter the local
neighborhood intersections, which in turn affect edge masses and conductances.
The implementation is not a simple direct weighted Laplacian with
\(w_{ij}=1/(\ell_{ij}+\epsilon)\).

\section{Methodological Caveat and Recommended Comparison}

The current \code{fit.rdgraph.regression()} operator should therefore be
described as an overlap-density/Riemannian-complex smoother. It is not the most
standard metric-length graph smoother. For the SIMODS geodesic-data-geometry
manuscript, this distinction matters because the non-oracle graph-selection
objective is meant to assess graph metric geometry. A more standard graph-signal
alternative would construct conductance directly from the supplied edge length,
for example
\[
  c_e^{\mathrm{len}}
    =
    \frac{1}{\ell_e+\epsilon}.
\]
That length-based conductance rule should be benchmarked as an alternative
smoother, not presented as the current behavior of
\code{fit.rdgraph.regression()}.

For future non-oracle graph-parameter selection experiments, we recommend a
controlled comparison between the current overlap-density conductance
\[
  c_e^\rho
    =
    \frac{1}{\max(\rho_1(e),10^{-10})}
\]
and a direct length-based conductance
\[
  c_e^{\mathrm{len}}
    =
    \frac{1}{\ell_e+\epsilon}.
\]
The key benchmark question is whether these two smoothers select graph
parameters that agree with oracle geodesic-isometry performance on synthetic
quad benchmarks. This comparison would separate two scientific questions that
are currently entangled: whether a candidate graph has good local overlap
geometry under the Riemannian-complex construction, and whether its metric edge
lengths support good graph-signal smoothing.

\section{Spectral Low-Pass Filter}

Let \((\lambda_j,v_j)\), \(j=1,\ldots,m\), be the cached eigenpairs of the
operator used for smoothing, normally \(\widetilde L_0\). Let
\[
  V = [v_1,\ldots,v_m],
  \qquad
  \Lambda=\diag(\lambda_1,\ldots,\lambda_m).
\]
For a fully observed response \(y\in\mathbb{R}^n\), the smoother evaluates
\[
  y_{\mathrm{spec}} = V^\top y,
  \qquad
  \hat y_\eta
  =
  V\, h_\eta(\Lambda)\, V^\top y,
\]
where \(h_\eta(\Lambda)\) is diagonal with entries \(h_\eta(\lambda_j)\).

The supported filter families are:
\[
\begin{array}{ll}
\texttt{heat\_kernel}:&
  h_\eta(\lambda)=\exp(-\eta\lambda),
\\[0.35em]
\texttt{tikhonov}:&
  h_\eta(\lambda)=\dfrac{1}{1+\eta\lambda},
\\[0.75em]
\texttt{cubic\_spline}:&
  h_\eta(\lambda)=\dfrac{1}{1+\eta\lambda^2},
\\[0.75em]
\texttt{gaussian}:&
  h_\eta(\lambda)=\exp(-\eta\lambda^2),
\\[0.35em]
\texttt{exponential}:&
  h_\eta(\lambda)=\exp(-\eta\sqrt{\max(\lambda,0)}),
\\[0.35em]
\texttt{butterworth}:&
  h_\eta(\lambda)=\dfrac{1}{1+(\lambda/\eta)^4}.
\end{array}
\]
The R wrapper default is \code{filter.type = "heat\_kernel"}.

\section{GCV Criterion and Eta Grid}

\subsection{GCV in \code{fit.rdgraph.regression()}}

For a fully labeled response, \code{fit.rdgraph.regression()} evaluates 40
candidate \(\eta\) values and selects the one minimizing
\[
  \operatorname{GCV}(\eta)
  =
  \frac{\|y-\hat y_\eta\|_2^2}
       {(n-\tr(S_\eta))^2},
  \qquad
  \tr(S_\eta)=\sum_{j=1}^m h_\eta(\lambda_j).
\]
The implementation floors the denominator at \(10^{-10}\) in absolute value.
For each candidate, it computes filter weights, applies them in the spectral
domain, reconstructs \(\hat y_\eta\), computes the GCV score, and retains the
best candidate.

The C++ eta grid in \code{fit.rdgraph.regression()} is logarithmic with
\(\eta_{\min}=10^{-10}\). Its \(\eta_{\max}\) depends on the filter family and
the largest retained eigenvalue \(\lambda_{\max}\):
\[
  \eta_{\max}
  =
  \begin{cases}
  -\log(10^{-10})/\lambda_{\max},&
    \text{heat kernel, Gaussian, exponential},\\
  10^{10},&
    \text{Tikhonov, cubic spline, Butterworth},\\
  1,&
    \text{fallback}.
  \end{cases}
\]
The grid is generated by exponential spacing between these bounds.

The fitter records the GCV result after the initial smoothing and after each
geometric refinement iteration. The returned fitted values are selected from
the iteration whose recorded GCV score is minimal. At the end, the fitter also
stores the filter weights \(h_{\eta^\star}(\lambda_j)\) for the selected
iteration.

\subsection{Semi-Supervised GCV}

When \code{y.vertices} is supplied to \code{fit.rdgraph.regression()}, the C++
smoother uses a masked spectral Tikhonov system rather than the direct
\(Vh_\eta(\Lambda)V^\top y\) formula. Let \(P\) be the diagonal mask for
labeled vertices \(L\), and let
\[
  A = V^\top P V,\qquad b=V^\top P y.
\]
For each \(\eta\), it solves
\[
  (A+\eta\Lambda)a=b,
  \qquad
  \hat y_\eta = Va,
\]
and uses the labeled-set GCV criterion
\[
  \operatorname{GCV}_L(\eta)
  =
  \frac{\|P(y-\hat y_\eta)\|_2^2}
       {(n_{\mathrm{eff}}-\mathrm{df}(\eta))^2},
  \qquad
  n_{\mathrm{eff}}=|L|,
\]
with
\[
  \mathrm{df}(\eta)
    =
    \tr\{(A+\eta\Lambda)^{-1}A\}.
\]

\section{Refit and Application to Multiple Features}

\code{refit.rdgraph.regression()} expects a fitted object of class
\code{knn.riem.fit}. It extracts:
\[
  V,\quad
  \lambda_1,\ldots,\lambda_m,\quad
  h_{\eta^\star}(\lambda_1),\ldots,h_{\eta^\star}(\lambda_m),\quad
  \eta^\star.
\]

With the default \code{per.column.gcv = FALSE} and no \code{y.vertices}, refit
applies the cached linear smoother to each column of a new response matrix
\(Y\):
\[
  \widehat Y
  =
  V\,\diag(h_{\eta^\star}(\lambda_1),\ldots,h_{\eta^\star}(\lambda_m))\,V^\top Y.
\]
The R implementation computes \(V^\top Y\), multiplies each spectral row by the
cached filter weight, and maps back by \(V\). For large matrices, it can do this
in column blocks.

If \code{per.column.gcv = TRUE}, refit does not reuse the original
\(\eta^\star\). Instead, it reuses the cached eigenbasis and eigenvalues, builds
an R-side eta grid, computes filter weights for each candidate, and selects an
independent \(\eta_j^\star\) for each response column by GCV. The R-side grid
uses finite positive eigenvalues, with
\[
  \eta_{\mathrm{lo}}=1/\max(\lambda_j),
  \qquad
  \eta_{\mathrm{hi}}=1/\min(\lambda_j),
\]
and logarithmic spacing; it falls back to a linear grid from \(10^{-3}\) to 1
if the eigenvalue range is unusable. Per-column GCV is disabled when
\code{y.vertices} is supplied.

If \code{y.vertices} is supplied to refit, the default fixed-\(\eta\) path
solves
\[
  (V_L^\top V_L+\eta^\star\Lambda)a=V_L^\top y_L,
  \qquad
  \hat y=Va,
\]
for each column or column block. This is a semi-supervised fixed-\(\eta\)
application, not a per-column GCV refit.

\section{Implications for Non-Oracle Graph Selection}

For graph selection, the recommended implementation-facing protocol is:
\begin{enumerate}[leftmargin=1.5em]
  \item For each candidate graph \(G_r\), pass \code{adj.list} and
  \code{weight.list} to \code{fit.rdgraph.regression()} with a proxy response or
  selected feature signal.
  \item Ensure \code{weight.list} contains positive symmetric edge lengths,
  not affinities or pre-inverted conductances.
  \item Extract the fitted model's GCV summary from \code{fit\$gcv}; the fitted
  values correspond to the iteration with minimum GCV.
  \item Apply \code{refit.rdgraph.regression()} to a matrix of selected feature
  signals. Use \code{per.column.gcv = TRUE} if the graph is to be scored by
  feature-specific smoothness parameters; otherwise use the cached original
  filter for all columns.
  \item Aggregate the graph score over selected features, for example by the
  mean or median of per-column GCV scores when \code{per.column.gcv = TRUE}.
\end{enumerate}

Manuscript-facing claims should describe this as a graph-spectral low-pass GCV
criterion based on the \code{gflow} Riemannian graph smoother. The manuscript
should not claim that edge lengths are directly transformed into affinities
\(1/(\ell+\epsilon)\). That formula is at best an intuition for inverse
edge-mass conductance after the neighborhood-intersection construction, not the
implemented edge-length conversion.

If the manuscript or benchmark wants a metric-length graph-smoothing selector,
that should be implemented and documented as a distinct comparator using a
conductance such as \(c_e^{\mathrm{len}}=1/(\ell_e+\epsilon)\). The current
\code{fit.rdgraph.regression()} selector should remain labeled as the
overlap-density/Riemannian-complex GCV smoother.

\section{Unresolved Ambiguities and Developer-Confirmation Items}

\begin{itemize}[leftmargin=1.5em]
  \item \textbf{Spectral operator versus full \(M_1\).} The generic Hodge
  operator \(L[0]\) uses the metric machinery, while the spectral smoother
  normally decomposes \code{L0\_mass\_sym}, built from diagonal conductances
  \(c_e=1/\rho_1(e)\). If off-diagonal \(M_1\) entries are enabled through
  triangle construction, those entries affect the generic \(L[0]\) but do not
  appear to enter \code{L0\_mass\_sym} except through the diagonal conductance
  vector. This appears intentional from the source comments, but should be
  treated as an implementation detail unless the gflow developer confirms it.
  \item \textbf{Benchmark scoring convention.} The exact aggregation rule over
  feature columns, and whether to use fixed-\(\eta\) refit or per-column GCV,
  is a benchmark-design decision, not fixed by the package.
  \item \textbf{Response-coherence modulation.} The fitter can iteratively
  update edge masses based on the response. For graph-selection benchmarks
  seeking a graph-only criterion, the experiment design should decide whether
  to keep \code{response.penalty.exp = 0} so the graph geometry remains fixed
  after initialization.
\end{itemize}

\section{Source Traceability}

\small
\begin{longtable}{p{0.25\linewidth}p{0.22\linewidth}p{0.31\linewidth}p{0.14\linewidth}}
\toprule
Mathematical or operational object & Code object & Source path and line region & Notes \\
\midrule
\endfirsthead
\toprule
Mathematical or operational object & Code object & Source path and line region & Notes \\
\midrule
\endhead
\bottomrule
\endfoot
R precomputed graph API & \code{adj.list}, \code{weight.list} arguments &
\Rpath{/Users/pgajer/current_projects/gflow/R/riem_dcx_regression.R}, lines 30--36, 849--950 &
\code{weight.list} documented as edge-length list; validates positivity, symmetry, and reciprocal agreement. \\

Default GCV graph-selection regime & function defaults &
\Rpath{/Users/pgajer/current_projects/gflow/R/riem_dcx_regression.R}, lines 654--681 &
Defaults include \code{response.penalty.exp = 0}, \code{use.counting.measure = TRUE}, \code{density.alpha = 1.5}, and \code{density.epsilon = 1e-10}. \\

Precomputed graph requires counting measure & R wrapper guard &
\Rpath{/Users/pgajer/current_projects/gflow/R/riem_dcx_regression.R}, lines 855--866 &
Supported R path stops if \code{use.counting.measure = FALSE}. \\

R-to-C++ call & \code{.Call(S\_fit\_rdgraph\_regression, ...)} &
\Rpath{/Users/pgajer/current_projects/gflow/R/riem_dcx_regression.R}, lines 1384--1425 &
Passes 0-based adjacency and numeric weights to C++. \\

Registered C entry point & \code{S\_fit\_rdgraph\_regression} &
\Rpath{/Users/pgajer/current_projects/gflow/src/init.c}, lines 329--332 &
Registers 42-argument C routine. \\

Precomputed graph extraction & \code{precomputed\_adj\_list}, \code{precomputed\_weight\_list} &
\Rpath{/Users/pgajer/current_projects/gflow/src/riem_dcx_regression_r.cpp}, lines 1429--1482 &
Copies weights but does not invert them. \\

Graph initialization dispatch & \code{fit\_rdgraph\_regression()} &
\Rpath{/Users/pgajer/current_projects/gflow/src/riem_dcx_regression.cpp}, lines 5924--6013 &
Calls \code{initialize\_from\_graph()} when precomputed graph is supplied. \\

Edge length use in neighborhoods & \code{initialize\_from\_graph()} &
\Rpath{/Users/pgajer/current_projects/gflow/src/riem_dcx_initialization.cpp}, lines 209--262, 515--558 &
Sorts by supplied lengths and builds pseudo-kNN neighborhoods. \\

Reference measure & \code{initialize\_reference\_measure()} &
\Rpath{/Users/pgajer/current_projects/gflow/src/riem_dcx_initialization.cpp}, lines 1490--1665 &
Counting measure is uniform; optional distance branch uses \((\epsilon+d_k)^{-\alpha}\). \\

Vertex and edge densities & \code{compute\_initial\_densities()} &
\Rpath{/Users/pgajer/current_projects/gflow/src/riem_dcx_regression.cpp}, lines 400--640 &
Defines \(\rho_0\) and \(\rho_1\) via neighborhood and intersection masses. \\

Mass matrices & \code{initialize\_metric\_from\_density()}, \code{compute\_edge\_mass\_matrix()} &
\Rpath{/Users/pgajer/current_projects/gflow/src/riem_dcx_regression.cpp}, lines 820--1068 &
Builds \(M_0\), \(M_1\), diagonal floors, and optional triple-intersection off-diagonals. \\

Boundary operator & \code{build\_boundary\_operator\_from\_edges()} &
\Rpath{/Users/pgajer/current_projects/gflow/src/riem_dcx_regression.cpp}, lines 1070--1126 &
Builds oriented incidence matrix \(B_1\). \\

Conductance convention & \code{assemble\_operators()} &
\Rpath{/Users/pgajer/current_projects/gflow/src/riem_dcx.cpp}, lines 289--307 &
Sets \(c_e=1/\max(\rho_1(e),10^{-10})\). \\

Mass-symmetrized Laplacian & \code{build\_L0\_mass\_sym\_if\_needed()} &
\Rpath{/Users/pgajer/current_projects/gflow/src/laplacian_bundle.cpp}, lines 15--188 &
Builds \(M_0^{-1/2}B_1CB_1^\top M_0^{-1/2}\). \\

Generic Hodge assembly & \code{laplacian\_bundle\_t::assemble\_all()} &
\Rpath{/Users/pgajer/current_projects/gflow/src/laplacian_bundle.cpp}, lines 190--245 &
Documents \(B_{p+1}M_{p+1}^{-1}B_{p+1}^\top M_p\) and lower-dimensional term. \\

Spectral decomposition target & \code{compute\_spectral\_decomposition()} &
\Rpath{/Users/pgajer/current_projects/gflow/src/riem_dcx_regression.cpp}, lines 1679--1698 &
Uses \code{L0\_mass\_sym} if available, otherwise \code{L[0]}. \\

Filter formulas & \code{apply\_filter\_function()} &
\Rpath{/Users/pgajer/current_projects/gflow/src/riem_dcx_regression.cpp}, lines 3990--4028 &
Defines heat, Tikhonov, cubic spline, Gaussian, exponential, and Butterworth filters. \\

Fit-time GCV and eta grid & \code{compute\_gcv\_score()}, \code{generate\_eta\_grid()} &
\Rpath{/Users/pgajer/current_projects/gflow/src/riem_dcx_regression.cpp}, lines 4050--4115 &
Defines RSS over \((n-\tr S)^2\) and C++ eta grid. \\

Fit-time smoother & \code{smooth\_response\_via\_spectral\_filter()} &
\Rpath{/Users/pgajer/current_projects/gflow/src/riem_dcx_regression.cpp}, lines 4123--4338 &
Computes eigenbasis projection, filter weights, GCV, and best \(\eta\). \\

Semi-supervised smoother & \code{smooth\_response\_via\_spectral\_filter\_semiy()} &
\Rpath{/Users/pgajer/current_projects/gflow/src/riem_dcx_regression.cpp}, lines 4360--4568 &
Uses masked Tikhonov system and labeled-set GCV. \\

Cached spectral component & \code{create\_spectral\_component()} &
\Rpath{/Users/pgajer/current_projects/gflow/src/riem_dcx_regression_r.cpp}, lines 500--640 &
Exports eigenvalues, filtered eigenvalues, eigenvectors, optimal eta, and filter type. \\

Fit result selected by GCV & result assembly &
\Rpath{/Users/pgajer/current_projects/gflow/src/riem_dcx_regression_r.cpp}, lines 2023--2060, 2266--2270 &
Selects fitted values from minimum-GCV iteration and returns spectral cache. \\

Fixed-\(\eta\) refit & \code{refit.rdgraph.regression()} &
\Rpath{/Users/pgajer/current_projects/gflow/R/refit_rdgraph_regression.R}, lines 111--115, 618--676 &
Applies cached \(V\), eigenvalues, \(\eta^\star\), and filter weights to new columns. \\

Per-column GCV refit & \code{compute.filter.weights.matrix()}, \code{select.eta.gcv.single()} &
\Rpath{/Users/pgajer/current_projects/gflow/R/refit_rdgraph_regression.R}, lines 270--306, 789--905 &
Reuses eigenbasis but selects per-column \(\eta\) with R-side GCV. \\

R-side eta grid for refit & \code{generate.eta.grid()} &
\Rpath{/Users/pgajer/current_projects/gflow/R/compat_helpers.R}, lines 86--104 &
Uses \(1/\lambda_{\max}\) to \(1/\lambda_{\min}\) log grid for per-column refit. \\

High-level data smoother & \code{data.smoother()} &
\Rpath{/Users/pgajer/current_projects/gflow/R/data_smoother.R}, lines 1--25, 400--486 &
Selects graph scale by fit-time GCV, then calls refit on the full matrix. \\

\end{longtable}
\normalsize

\end{document}
